% =====================================================================
%  Semantics frames (background deck)
%
%  B2 "Probabilistic semantics": the flip program  ≡  its semantics type
%       M : ⟦Γ⟧ → Measure ty
%  B3 "Refining the domain": the SAME type, same on-screen position, but
%     the domain grows by prepending  HC →  giving
%       HC (Hilbert cube) → ⟦Γ⟧ → Measure ty
%
%  DESIGN CHOICE -------------------------------------------------------
%  The "grows into place" effect is realised as a SINGLE frame with beamer
%  overlays (B3 is a continuation of B2 in the same frame). This guarantees
%  the type sits in EXACTLY the same position across overlays: the fixed
%  part  ⟦Γ⟧ → Measure ty  is typeset once and never moves; the  HC →
%  prefix is revealed on the second overlay and fades in via a tikz
%  fading. We anchor the whole expression to the RIGHT of a fixed point so
%  that prepending the prefix grows the formula leftwards into reserved
%  whitespace, leaving the tail (the part the audience tracks) stationary.
%
%  Rather than two separate frames, the single-frame-with-overlays form
%  keeps the layout provably identical (no risk of a 1pt drift between two
%  hand-laid frames) — cleanest continuous animation.
%
%  Macros from main.tex preamble: \sem{·} (⟦·⟧), \alert{}, colour words.
% =====================================================================

% ---------------------------------------------------------------------
%  Frame B2 / B3 — combined, overlay-driven
% ---------------------------------------------------------------------
\begin{frame}{%
    \only<1>{Probabilistic semantics}%
    \only<2->{Refining the domain: the Hilbert cube}%
  }

  \begin{center}
    \begin{minipage}[c]{0.40\textwidth}
      % --- LEFT: the small program, set in plain math (NOT minted) ------
      \centering
      \(
        \begin{aligned}
          X_0 &\leftarrow \mathsf{flip}\ 0\\
          X_1 &\leftarrow \mathsf{flip}\ X_0\\
          X_2 &\leftarrow \mathsf{flip}\ X_1\\
              &\mathsf{ret}\ X_2
        \end{aligned}
      \)
    \end{minipage}%
    \hfill
    \begin{minipage}[c]{0.06\textwidth}
      \centering\Large\(\equiv\)
    \end{minipage}%
    \hfill
    \begin{minipage}[c]{0.52\textwidth}
      % --- RIGHT: the semantics TYPE, large and prominent ---------------
      % The expression is right-anchored: the fixed tail never moves; the
      % HC → prefix grows leftwards into reserved whitespace.
      %
      % KEY trick: the prefix "HC →" RESERVES its full width on EVERY
      % overlay via a \phantom, so the fixed tail (⟦Γ⟧ → Measure ty) never
      % moves. The visible glyphs are overprinted in the same zero-extra-
      % width box (\rlap) and revealed only on overlay 2+, so the prefix
      % "grows into place" while the tail stays pinned in position.
      \raggedleft
      \Large
      \(
        M : {}
        \rlap{\(\only<2->{\alert{\mathrm{HC}}\;\alert{\to}}\)}%
        \phantom{\mathrm{HC}\;\to}\;
        % --- the FIXED tail: typeset once, identical on every overlay ----
        \sem{\Gamma} \to \mathrm{Measure}\ \mathit{ty}
      \)
      % --- small annotation under the HC prefix (overlay 2+) ------------
      \\[0.2em]
      \only<2->{\makebox[\linewidth][l]{\hspace*{0.2cm}%
        \alert{\footnotesize \(\mathrm{HC}\) = Hilbert cube \(=[0,1]^{\mathbb{N}}\)}}}
    \end{minipage}
  \end{center}

  \vfill

  % -------------------------------------------------------------------
  %  Stylised Hilbert-cube illustration (overlay 2+).
  %  A row of ~6 coordinate cells: scattered sample dots, a couple
  %  crossed out, thin curved arrows into a {X₁, X₂, X₃} label.
  % -------------------------------------------------------------------
  % stylised placeholder — swap for a nicer asset later
  \only<2->{%
  \begin{center}
  \begin{tikzpicture}[
      cell/.style={draw, gray, thin, minimum size=0.75cm, inner sep=0pt},
      dot/.style={circle, fill=wordblue, inner sep=0pt, minimum size=2.4pt},
      arr/.style={-{Stealth[length=1.6mm]}, thin, gray},
    ]
    % --- the coordinate cells -----------------------------------------
    \foreach \i in {0,...,5} {
      \node[cell] (c\i) at (\i*0.95,0) {};
    }
    % --- scattered sample dots inside a few cells ---------------------
    \node[dot] at ([shift={(-0.18, 0.20)}]c0.center) {};
    \node[dot] at ([shift={( 0.15,-0.10)}]c0.center) {};
    \node[dot] at ([shift={( 0.05, 0.25)}]c1.center) {};
    \node[dot] at ([shift={(-0.20,-0.18)}]c1.center) {};
    \node[dot] at ([shift={( 0.18, 0.12)}]c1.center) {};
    \node[dot] at ([shift={( 0.00, 0.00)}]c3.center) {};
    \node[dot] at ([shift={(-0.15, 0.18)}]c3.center) {};
    \node[dot] at ([shift={( 0.20,-0.20)}]c4.center) {};
    % --- a couple of cells crossed out (✗) ----------------------------
    \draw[red, thin] (c2.south west) -- (c2.north east);
    \draw[red, thin] (c2.north west) -- (c2.south east);
    \draw[red, thin] (c5.south west) -- (c5.north east);
    \draw[red, thin] (c5.north west) -- (c5.south east);
    % --- label to the right -------------------------------------------
    \node[anchor=west, font=\small] (lbl) at (6.4,0)
      {\(\{X_1, X_2, X_3\}\)};
    % --- thin arrows curving from active cells into the label ---------
    \draw[arr] (c0.north) to[out=60, in=160] (lbl.north west);
    \draw[arr] (c1.north) to[out=70, in=170] ([yshift=1mm]lbl.west);
    \draw[arr] (c4.south) to[out=-30, in=200] (lbl.south west);
  \end{tikzpicture}
  \end{center}%
  }
\end{frame}
